
% =============================================================================
% Project : THE PHILOSOPHY OF VIABILITY
% Book    : Book I — Ontology of Viability
% Book ID : VS.B01
% Chapter : Chapter 08 — Comparative Philosophical Examination
% Section : Section 14 — Suhrawardi
% Asset   : Diagram 01 — Orders of Light and Classes of Existent Things
% Asset ID: VS.B01.P01.C08.S14.SS09.P004.DG01
% File    : 4_ASSETS/diagrams/VS.B01.P01.C08.S14.SS09.P004.DG01.suhrawardi_orders_of_light.tex
%
% Qualification:
%   The two panels represent distinct classificatory dimensions.
%   The four classes on the right are not arranged as a causal sequence.
% =============================================================================

\begin{figure}[htbp]
\centering

\begin{tikzpicture}[
    >=Stealth,
    font=\small,
    lightnode/.style={
        draw=ViabilityDarkBlue,
        fill=ViabilityPaleBlue,
        rounded corners=2pt,
        align=center,
        text width=5.0cm,
        minimum height=10mm,
        inner sep=5pt
    },
    ordernode/.style={
        draw=ViabilityGold,
        fill=ViabilityPaleGold,
        rounded corners=2pt,
        align=center,
        text width=4.8cm,
        minimum height=12mm,
        inner sep=5pt
    },
    classnode/.style={
        draw=ViabilityGreen,
        fill=ViabilityPaleGreen,
        rounded corners=2pt,
        align=center,
        text width=5.0cm,
        minimum height=11mm,
        inner sep=5pt
    },
    paneltitle/.style={
        font=\bfseries,
        align=center,
        text width=5.2cm
    }
]

% Panel A: Orders of lights

\node[paneltitle] at (0,0)
    {Two Orders of Immaterial Lights};

\node[lightnode] (lol) at (0,-1.0)
    {\textbf{Light of Lights}\\
     First principle of the Illuminationist order};

\node[ordernode] (vertical) at (-1.45,-2.7)
    {\textbf{Vertical Order}\\
     Dominating lights and ontological dependence};

\node[ordernode] (horizontal) at (1.45,-2.7)
    {\textbf{Horizontal Order}\\
     Equal-rank lights, mother-lights, and Lords of Species};

\node[lightnode] (managing) at (0,-4.45)
    {\textbf{Managing Lights}\\
     Soul-related lights associated with corporeal orders};

\draw[->,thick,draw=ViabilityDarkBlue]
    (lol.south) -- (vertical.north);

\draw[<->,dashed,thick,draw=ViabilityGold]
    (vertical.east) -- (horizontal.west);

\draw[densely dotted,thick,draw=ViabilityDarkBlue]
    (vertical.south) -- (managing.north);

\node[align=center,font=\scriptsize,text width=5.2cm]
    at (0,-5.35)
    {The dotted connection indicates association, not a further
     stage in a single linear sequence};

% Panel B: Four classes

\node[paneltitle] at (8.2,0)
    {Four Principal Classes of Existent Things};

\node[classnode] at (8.2,-1.0)
    {\textbf{Pure Immaterial Lights}\\
     Self-subsistent and incorporeal};

\node[classnode] at (8.2,-2.2)
    {\textbf{Dark Barriers}\\
     Corporeal and extended entities};

\node[classnode] at (8.2,-3.4)
    {\textbf{Accidental Lights}\\
     Luminous states inhering in another};

\node[classnode] at (8.2,-4.6)
    {\textbf{Dark Accidents}\\
     Non-luminous states inhering in another};

\node[
    draw=gray,
    fill=white,
    rounded corners=2pt,
    align=center,
    text width=11.0cm,
    inner sep=5pt
] at (4.1,-6.0)
    {\textit{The two panels show different classificatory dimensions.
    No one-to-one correspondence between them is implied.}};

\end{tikzpicture}

\caption{Schematic distinction between Suhrawardi's vertical and
horizontal orders of immaterial lights and his four principal
classes of existent things. The diagram is non-temporal and
does not imply that the four classes form a causal sequence.}
\label{fig:ch08-s14-suhrawardi-orders-of-light}

\end{figure}

% =============================================================================
% END OF DIAGRAM 01 — SUHRAWARDI
% =============================================================================
